\documentclass[12pt]{article}
\usepackage[T1]{fontenc}
\usepackage[utf8]{inputenc}
\usepackage{lmodern}
\usepackage{textcomp}
\usepackage{enumitem}
\usepackage{geometry}
\geometry{margin=1in}
\begin{document}
\begin{center}
Answer Key - Science - K-12 Level Assessment 8 \
\small (Answers are ordered exactly as in the question paper.)
\end{center}

\section*{Multiple Choice}
\begin{enumerate}[label=\arabic*.]
\item \textbf{Answer:} D
\\ \textit{Options:} A) glass; B) metal; C) water; D) wood
\\ \textit{Note:} Wood is considered the best heat insulator because it has a low thermal conductivity, which means it does not easily allow heat to pass through it, making it effective at reducing heat transfer.
\item \textbf{Answer:} A
\\ \textit{Options:} A) specialize.; B) reproduce.; C) use energy.; D) make protein.
\\ \textit{Note:} The correct answer is "specialize" because single-cell organisms perform all necessary life functions independently, while cells in multicellular organisms can differentiate into specialized types to perform specific roles, contributing to the overall function of the organism.
\item \textbf{Answer:} D
\\ \textit{Options:} A) 0; B) 1; C) 4; D) 7
\\ \textit{Note:} All planets in our solar system, including the seven other planets besides Earth, orbit the Sun in the same direction, which is counterclockwise when viewed from above the Sun's north pole.
\item \textbf{Answer:} B
\\ \textit{Options:} A) valleys carved by a moving glacier; B) piles of rocks deposited by a melting glacier; C) grooves created in a granite surface by a glacier; D) bedrock hills roughened by the passing of a glacier
\\ \textit{Note:} The correct answer highlights that glaciers, as constructive forces, shape the landscape by depositing debris and rocks, known as moraines, as they melt and retreat, creating distinct landforms.
\item \textbf{Answer:} B
\\ \textit{Options:} A) 15^\ensuremath{\circC}; B) 21^\ensuremath{\circC}; C) 27^\ensuremath{\circC}; D) 33^\ensuremath{\circC}
\\ \textit{Note:} The median temperature is calculated by arranging the temperatures in ascending order (12, 18, 24, 31) and finding the average of the two middle values (18 and 24), which results in a median of 21 degrees Celsius.
\end{enumerate}

\section*{True / False}
\begin{enumerate}[label=\arabic*.]
\setcounter{enumi}{5}
\item \textbf{Answer:} True
\\ \textit{Options:} A) True; B) False
\\ \textit{Note:} Newts have the ability to autotomize, or shed their tails, as a defense mechanism to distract predators and increase their chances of survival.
\item \textbf{Answer:} False
\\ \textit{Options:} A) True; B) False
\\ \textit{Note:} The marketing department is primarily responsible for promoting and selling products, while the logistics or supply chain department typically manages shipping and storage.
\item \textbf{Answer:} True
\\ \textit{Options:} A) True; B) False
\\ \textit{Note:} Roasting marshmallows over a fire causes a chemical change because the heat breaks down the sugars in the marshmallow, resulting in the formation of new substances such as caramelized sugars and other compounds.
\item \textbf{Answer:} False
\\ \textit{Options:} A) True; B) False
\\ \textit{Note:} Hydraulic systems are closed systems that use incompressible fluids to transmit force, while pneumatic systems are open systems that use compressible gases to transmit force.
\item \textbf{Answer:} False
\\ \textit{Options:} A) True; B) False
\\ \textit{Note:} The statement is false because the Sun is much larger than Jupiter, making the correct order from largest to smallest Sun, Jupiter, Earth, Moon.
\end{enumerate}

\section*{Long Form Response}
\begin{enumerate}[label=\arabic*.]
\setcounter{enumi}{10}
\item \textbf{Answer:} To convert a height measurement of 6 feet into metric units, you would first need to know the conversion factor; 1 foot is equal to 0.3048 meters. Therefore, you multiply 6 feet by 0.3048 to get the height in meters: 6 x 0.3048 = 1.8288 meters. This process demonstrates the importance of using different measurement systems in science, as it allows for consistency and clarity when sharing data across various regions and scientific fields. Using metric units is often preferred in scientific contexts because they are based on multiples of ten, making calculations easier and more straightforward. Understanding how to convert measurements ensures accurate communication of scientific information globally.
\\ \textit{Note:} correct because it accurately applies the conversion factor from feet to meters and emphasizes the significance of using a standardized measurement system for clarity and consistency in scientific communication.
\item \textbf{Answer:} Geological processes such as tectonic activity, faulting, and sedimentation can lead to the arrangement of rock layers where younger layers are located beneath older layers, a situation known as an inverted stratigraphy. This can occur when geological forces push older rocks upwards while younger rocks settle down in depressions or basins. Understanding these arrangements is significant because they provide insights into Earth's history, indicating past environmental conditions, the sequence of geological events, and the age of rocks through relative dating. By studying these layers, scientists can reconstruct Earth's geological timeline and better understand the processes that have shaped our planet.
\\ \textit{Note:} The answer correctly identifies that geological processes such as tectonic activity and sedimentation can lead to inverted stratigraphy, illustrating how these arrangements reveal important insights into Earth's geological history and past environmental conditions.
\end{enumerate}

\vfill
\noindent\textit{End of Answer Key}
\end{document}